\scanpage{224}
que \deleted{\unsure{par}} sur
\[
\boxed{(6)\quad \underline\Phi=f_*(Z\underline\Omega^1_{X/S})/df_*(\underline{\mathcal O}_X)},
\]
il y a une connexion intégrable naturelle, grâce à la formule de Cartan \unsure{---|.} les $\omega_i$ définissent des sections $\omega'_i$ de $\underline\Phi$, et l'on peut former les $D_i\deleted{\unsure{\omega}}\omega'_i$. Je dis qu'il suffit qu'on ait
\margin{attention, \unsure{ne} marche tel quel que si $X\longrightarrow S$ est steinien}
\[
(7)\qquad \sum D_i\omega'_i=0.
\]

Lorsque $X/S$ est muni d'une connexion intégrable, alors le calcul des $\sum D_i\omega'_i$ se fait ainsi : \inserted{$f_*(\underline\Omega^*_{X/S})$ sera aussi stratifié; donc} chaque $D_i$ se remonte \deleted{par la stratification en} \inserted{en un opérateur} $\overline D_i$ sur $\underline\Omega^*_X$ \inserted{commutant à la diff. rel. $d_{X/S}$} et \deleted{\unsure{comme}} \deleted{$\omega_i$ est une forme fermée} \deleted{on peut former}

\deleted{\[
\boxed{\begin{gathered}
(8)\qquad \alpha_{\overline D_i}\omega_i=\alpha_{\overline D_i^0}\omega_i+c_i\omega_i\\
\left\{\begin{aligned}
\overline D_i&=\overline D_i^0+c_i\,\mathrm{id}\qquad
\left(\begin{array}{l}\overline D_i^0\text{ à coeff. cts}\\c_i\text{ section de }\mathcal O_X\end{array}\right)\\
(9)\qquad \alpha_{\overline D_i^0}d\varphi&=\unsure{d}\overline D_i^0\varphi
\end{aligned}\right.
\end{gathered}}
\]}

Considérons alors
\[
(\unsure{8})\qquad \sum\deleted{\unsure{\alpha}}(\overline D_i)\omega_i=\overline\omega
\]
qui est encore une forme fermée.
On vérifie alors \inserted{aussitôt} que l'on a
\margin{223}
